\section*{RESULTS}

\subsection{Circuit architecture: E/I modules as the oscillatory core}

How does a network of non-bursting neurons generate robust rhythmic output?
The \textit{Dendronotus} swim CPG achieves this through network hysteresis: slow synaptic feedback within contralateral E/I modules creates the history-dependence that individual neurons lack, driving alternation between left and right activity states (Fig.~\ref{fig:ei_mechanism}).
We constructed a computational model of this four-cell network by integrating validated two-cell subnetworks---half-center oscillators (HCOs) and E/I modules---following the approach developed in \cite{scully2024pairing}.

The circuit architecture distributes rhythmogenic function across the network rather than concentrating it in pacemaker cells (Fig.~\ref{fig:ei_mechanism}A).
Each E/I module pairs an excitatory Si3 neuron with a contralateral Si2 neuron in a ``twisted'' configuration that crosses the midline: Si3R drives Si2L and Si3L drives Si2R.
Within each module, Si3 provides excitatory drive through a logistic synapse with relatively fast kinetics, while Si2 returns slower inhibitory feedback through a logistic synapse.
This delayed negative feedback loop constitutes the conceptual and computational primitive of the swim network---the fundamental building block from which network-level bursting emerges.

Critically, each E/I module is tuned to be ``almost-oscillatory'': when isolated, it exhibits damped oscillations or subthreshold rhythmic fluctuations rather than sustained limit cycles.
This design means that no single component can generate rhythm autonomously; network coupling is required to recruit the modules into coordinated bursting.
The two HCOs (Si3L$\leftrightarrow$Si3R and Si2L$\leftrightarrow$Si2R) provide reciprocal inhibition that enforces anti-phase alternation between left and right sides.
Unlike classical HCO-based CPGs where reciprocal inhibition drives rhythm generation, these connections function primarily as stabilizers---the E/I modules supply the oscillatory drive.
Weak electrical coupling between E/I partners (Si3R$\leftrightarrow$Si2L and Si3L$\leftrightarrow$Si2R) provides rapid voltage coordination within each module and contributes to network resilience.

The kinetic asymmetry between fast excitation and slow inhibition is visible in the neurotransmitter traces (Fig.~\ref{fig:ei_mechanism}B): excitatory synaptic variables (red) rise and fall rapidly with each spike, while inhibitory variables (blue) accumulate gradually across the burst and decay slowly during interburst intervals.
This temporal separation means that Si3 excitation recruits Si2 activity quickly, but the inhibition that eventually terminates Si3 firing builds over hundreds of milliseconds---the timescale that sets burst duration.

These synaptic timescales are not arbitrary modeling choices but reflect dynamics observed in the biological circuit.
Figure~\ref{fig:synaptic_timescales} demonstrates correspondence between model and experimental recordings: current injection into Si2R evokes dual-timescale responses in its network partners.
The fast $\alpha$-synapse produces IPSPs in Si2L that rise and decay within tens of milliseconds.
Electrical coupling appears as an immediate voltage deflection in Si3L that precedes the slower chemical response.
The logistic synapse from Si2 to Si3 accumulates over repeated stimulation, producing delayed inhibition that builds across hundreds of milliseconds.
Together, fast reciprocal inhibition, slow E/I feedback, and near-instantaneous electrical coupling establish the temporal framework within which network bursting emerges.
